\chapter{Pigeon-hole and double counting}
\label{chapter28}

Some mathematical principles, such as the two in the title of this chapter,
are so obvious that you might think they would only produce equally
obvious results. To convince you that ``It ain't necessarily so'' we
illustrate them with examples that were suggested by Paul Erd\H{o}s to be
included in The Book. We will encounter instances of them also in later
chapters.

\begin{theorem}[Pigeon-hole principle]
  \label{pigeon_hole_principle}
  \lean{chapter28.pigeon_hole_principle}
  If $n$ objects are placed in $r$ boxes, where $r < n$, then at least one of the boxes contains
  more than one object.
\end{theorem}

Well, this is indeed obvious, there is nothing to prove. In the language of
mappings our principle reads as follows: Let $N$ and $R$ be two finite sets
with $|N| = n > r = |R|$, and let $f : N \to R$ be a mapping. Then there
exists some $a \in R$ with $|f^{-1}(a)| \geq 2$. We may even state a stronger
inequality: There exists some $a \in R$ with
\[
|f^{-1}(a)| \geq \left\lceil \frac{n}{r} \right\rceil. \tag{1}
\]
In fact, otherwise we would have $|f^{-1}(a)| < \frac{n}{r}$ for all $a$, and hence
$n = \sum_{a \in R} |f^{-1}(a)| < r \cdot \frac{n}{r} = n$, which cannot be.

\begin{proof}
  \leanok
  Obvious.
\end{proof}

\section{Numbers}

\begin{theorem}[Claim]
  \label{ch28claim1}
  \lean{chapter28.claim1_coprime}
  Consider the numbers $1, 2, 3, \dots, 2n$, and take any $n + 1$
  of them. Then there are two among these $n + 1$ numbers which are
  relatively prime.
\end{theorem}
\begin{proof}
  \leanok
  \uses{pigeon_hole_principle}
  This is again obvious. There must be two numbers which are only $1$ apart,
  and hence relatively prime.
\end{proof}

But let us now turn the condition around.

\begin{theorem}[Claim]
  \label{ch28claim2}
  \lean{chapter28.claim2_divisible}
  Suppose again $A \subset \{1, 2, \dots, 2n\}$ with $|A| = n+1$. Then
  there are always two numbers in $A$ such that one divides the other.
\end{theorem}
\begin{proof}
  \leanok
  \uses{pigeon_hole_principle}
  Write every number $a \in A$ in the form $a = 2^k m$, where $m$ is an odd number between $1$
  and $2n - 1$. Since there are $n + 1$ numbers in $A$, but only $n$ different odd parts, there must
  be two numbers in $A$ with the same odd part. Hence one is a multiple of the other.
\end{proof}

\section{Sequences}

\begin{theorem}[Claim (Erd\H{o}s--Szekeres)]
  \label{ch28claim3}
  \lean{chapter28.claim3_erdos_szekeres}
  In any sequence $a_1, a_2, \ldots, a_{mn+1}$ of $mn+1$ distinct real numbers,
  there exists an increasing subsequence
  \[
  a_{i_1} < a_{i_2} < \cdots < a_{i_{m+1}} \quad (i_1 < i_2 < \cdots < i_{m+1})
  \]
  of length $m+1$, or a decreasing subsequence
  \[
  a_{j_1} > a_{j_2} > \cdots > a_{j_{n+1}} \quad (j_1 < j_2 < \cdots < j_{n+1})
  \]
  of length $n+1$, or both.
\end{theorem}
\begin{proof}
  \leanok
  \uses{pigeon_hole_principle}
  This time the application of the pigeon-hole principle is not immediate.
  Associate to each $a_i$ the number $t_i$, which is the length of a longest
  increasing subsequence starting at $a_i$. If $t_i \geq m+1$ for some $i$,
  then we have an increasing subsequence of length $m+1$. Suppose then
  that $t_i \leq m$ for all $i$. The function $f : a_i \mapsto t_i$ mapping
  $\{a_1, \dots, a_{mn+1}\}$ to $\{1, \dots, m\}$ tells us by~(1) that there
  is some $s \in \{1, \dots, m\}$ such that $f(a_i) = s$ for $\frac{mn}{m} + 1 = n + 1$
  numbers $a_i$. Let $a_{j_1}, a_{j_2}, \dots, a_{j_{n+1}}$ $(j_1 < \cdots < j_{n+1})$
  be these numbers. Now look at two consecutive numbers $a_{j_i}, a_{j_{i+1}}$. If
  $a_{j_i} < a_{j_{i+1}}$, then we would obtain an increasing subsequence of length
  $s$ starting at $a_{j_{i+1}}$, and consequently an increasing subsequence of length
  $s+1$ starting at $a_{j_i}$, which cannot be since $f(a_{j_i}) = s$.
  We thus obtain a decreasing subsequence $a_{j_1} > a_{j_2} > \cdots > a_{j_{n+1}}$
  of length $n+1$.
\end{proof}

\section{Sums}

\begin{theorem}[Claim]
  \label{ch28claim4}
  \lean{chapter28.claim4_contiguous_sum}
  Suppose we are given $n$ integers $a_1, \dots, a_n$, which need not be distinct.
  Then there is always a set of consecutive numbers $a_{k+1}, a_{k+2}, \dots, a_{\ell}$
  whose sum $\sum_{i=k+1}^{\ell} a_i$ is a multiple of $n$.
\end{theorem}
\begin{proof}
  \leanok
  \uses{pigeon_hole_principle}
  For the proof we set $N = \{0, 1, \dots, n\}$ and $R = \{0, 1, \dots, n-1\}$.
  Consider the map $f : N \to R$, where $f(m)$ is the remainder of $a_1 + \cdots + a_m$
  upon division by $n$. Since $|N| = n+1 > n = |R|$,
  it follows that there are two sums $a_1 + \cdots + a_k, a_1 + \cdots + a_{\ell}$ $(k < \ell)$
  with the same remainder, where the first sum may be the empty sum denoted by $0$. It follows that
  \[
  \sum_{i=k+1}^{\ell} a_i = \sum_{i=1}^{\ell} a_i - \sum_{i=1}^{k} a_i
  \]
  has remainder $0$ --- end of proof.
\end{proof}

\bigskip

Let us turn to the second principle: counting in two ways. By this we mean
the following.

\begin{theorem}[Double counting]
  \label{double_counting}
  \lean{chapter28.double_counting}
  Suppose that we are given two finite sets $R$ and $C$ and a subset $S \subseteq R \times C$.
  Whenever $(p, q) \in S$, then we say $p$ and $q$ are incident.
  If $r_p$ denotes the number of elements that are incident to $p \in R$, and $c_q$
  denotes the number of elements that are incident to $q \in C$, then
  \[
  \sum_{p \in R} r_p = |S| = \sum_{q \in C} c_q. \tag{3}
  \]
\end{theorem}
\begin{proof}
  \leanok
  Again, there is nothing to prove. The first sum classifies the pairs in $S$
  according to the first entry, while the second sum classifies the same pairs
  according to the second entry.
\end{proof}

There is a useful way to picture the set $S$. Consider the matrix $A = (a_{pq})$,
the incidence matrix of $S$, where the rows and columns of $A$ are indexed
by the elements of $R$ and $C$, respectively, with
\[
a_{pq} = \begin{cases} 1 & \text{if } (p, q) \in S, \\ 0 & \text{if } (p, q) \notin S. \end{cases}
\]
With this set-up, $r_p$ is the sum of the $p$-th row of $A$ and $c_q$ is the sum of the
$q$-th column. Hence the first sum in~(3) adds the entries of $A$ (that is, counts
the elements in $S$) by rows, and the second sum by columns.

\section{Numbers again}

Look at the incidence matrix for $R = C = \{1, 2, \dots, n\}$ and $S = \{(i,j) : i \mid j\}$.
The number of $1$'s in column $j$ is precisely the number of divisors of $j$; let us
denote this number by $t(j)$. Let us ask how large this number $t(j)$ is on the average
when $j$ ranges from $1$ to $n$. Thus, we ask for the quantity
\[
\bar{t}(n) = \frac{1}{n} \sum_{j=1}^{n} t(j).
\]

How large is $\bar{t}(n)$ for arbitrary $n$? At first glance, this seems hopeless. For
prime numbers $p$ we have $t(p) = 2$, while for $2^k$ we obtain a large number
$t(2^k) = k + 1$. So, $t(n)$ is a wildly jumping function, and we surmise that
the same is true for $\bar{t}(n)$. Wrong guess, the opposite is true! Counting in
two ways provides an unexpected and simple answer.

\begin{theorem}[Average number of divisors]
  \label{ch28_avg_divisors}
  \lean{chapter28.sum_divisor_count}
  For any positive integer $n$,
  \[
  \sum_{j=1}^{n} t(j) = \sum_{i=1}^{n} \left\lfloor \frac{n}{i} \right\rfloor.
  \]
  Consequently,
  \[
  \log n - 1 < \bar{t}(n) \leq H_n < \log n + 1,
  \]
  where $H_n = \sum_{i=1}^{n} \frac{1}{i}$ is the $n$-th harmonic number.
\end{theorem}
\begin{proof}
  \leanok
  \uses{double_counting}
  Consider the matrix $A$ (as above) for the integers $1$ up to $n$. Counting by
  columns we get $\sum_{j=1}^{n} t(j)$. How many $1$'s are in row $i$? Easy enough, the
  $1$'s correspond to the multiples of $i$: $1 \cdot i, 2 \cdot i, \dots$, and the last multiple not
  exceeding $n$ is $\lfloor n/i \rfloor \cdot i$. Hence we obtain
  \[
  \bar{t}(n) = \frac{1}{n} \sum_{j=1}^{n} t(j) = \frac{1}{n} \sum_{i=1}^{n} \left\lfloor \frac{n}{i} \right\rfloor
  \leq \frac{1}{n} \sum_{i=1}^{n} \frac{n}{i} = \sum_{i=1}^{n} \frac{1}{i} = H_n,
  \]
  where the error in each summand, when passing from $\lfloor n/i \rfloor$ to $n/i$, is less
  than $1$. Together with the standard estimates $\log n < H_n < \log n + 1$
  this gives $\log n - 1 < \bar{t}(n) \leq H_n < \log n + 1$.

  Thus we have proved the remarkable result that, while $t(n)$ is totally erratic,
  the average $\bar{t}(n)$ behaves beautifully: It differs from $\log n$ by less than $1$.
\end{proof}

\section{Graphs}

Let $G$ be a finite simple graph with vertex set $V$ and edge set $E$. The
degree $d(v)$ of a vertex $v$ is the number of edges which have $v$ as an
end-vertex.

Almost every book in graph theory starts with the following result:

\begin{lemma}[Handshaking]
  \label{handshaking}
  \lean{chapter28.handshaking}
  \leanok
  Let $G$ be a finite simple graph with vertex set $V$ and edge set $E$ and let
  $d(v)$ denote the degree of a vertex $v$. Then
  \[
  \sum_{v\in V}d(v) = 2|E|. \tag{4}
  \]
\end{lemma}
\begin{proof}
  \leanok
  For the proof consider $S \subseteq V \times E$, where $S$ is
  the set of pairs $(v, e)$ such that $v \in V$ is an end-vertex
  of $e \in E$. Counting $S$ in two ways gives on the one hand
  $\sum_{v \in V} d(v)$, since every vertex contributes $d(v)$
  to the count, and on the other hand $2|E|$, since every edge
  has two ends.
\end{proof}

We want to single out the following beautiful application to an extremal problem
on graphs. Here is the problem:
\begin{quote}
Suppose $G = (V, E)$ has $n$ vertices and contains no cycle of length $4$ (denoted by $C_4$).
How many edges can $G$ have at most?
\end{quote}

Let us tackle the general problem. Let $G$ be a graph on $n$ vertices without
a $4$-cycle. We count the following set $S$ in two ways: $S$ is the set of pairs
$(u, \{v, w\})$ where $u$ is adjacent to $v$ and to $w$, with $v \neq w$.
In other words, we count all occurrences of a ``cherry'' (path of length~$2$ centered at~$u$).

Summing over $u$, we find $|S| = \sum_{u \in V} \binom{d(u)}{2}$. On the other hand,
every pair $\{v, w\}$ has at most one common neighbor (by the $C_4$-condition).
Hence $|S| \leq \binom{n}{2}$, and we conclude
\[
\sum_{u \in V} \binom{d(u)}{2} \leq \binom{n}{2}
\]
or
\[
\sum_{u \in V} d(u)^2 \leq n(n-1) + \sum_{u \in V} d(u). \tag{5}
\]

Next we apply the Cauchy--Schwarz inequality to the vectors $(d(u_1), \dots, d(u_n))$
and $(1, 1, \dots, 1)$, obtaining
\[
\left(\sum_{u \in V} d(u)\right)^2 \leq n \sum_{u \in V} d(u)^2,
\]
and hence by~(5)
\[
\left(\sum_{u \in V} d(u)\right)^2 \leq n^2(n-1) + n \sum_{u \in V} d(u).
\]
Invoking~(4) we find
\[
4|E|^2 \leq n^2(n-1) + 2n|E|
\]
or
\[
|E|^2 - \frac{n}{2}|E| - \frac{n^2(n-1)}{4} \leq 0.
\]
Solving the corresponding quadratic equation we thus obtain the following
result of Istv\'an Reiman.

\begin{theorem}[Reiman]
  \label{ch28theorem}
  \lean{chapter28.c4_free_edge_bound}
  If the graph $G$ on $n$ vertices contains no $4$-cycles, then
  \[
  |E| \leq \left\lfloor \frac{n}{4}\left(1 + \sqrt{4n - 3}\right) \right\rfloor. \tag{6}
  \]
\end{theorem}
\begin{proof}
  \leanok
  \uses{double_counting, handshaking}
  The proof follows from the chain of inequalities above: the key combinatorial
  step $\sum \binom{d(u)}{2} \leq \binom{n}{2}$ is established by double counting
  (Theorem~\ref{ch28_sum_choose}), and the final bound follows by
  Cauchy--Schwarz and the quadratic formula.
\end{proof}

\begin{theorem}[Cherry counting]
  \label{ch28_sum_choose}
  \lean{chapter28.sum_choose_deg_le_choose_card}
  \leanok
  If the graph $G$ on $n$ vertices contains no $4$-cycles, then
  \[
  \sum_{v \in V} \binom{d(v)}{2} \leq \binom{n}{2}.
  \]
\end{theorem}
\begin{proof}
  \leanok
  \uses{double_counting}
  Each pair $(v, \{u, w\})$ with $u, w \in N(v)$ maps to $\{u, w\}$. The $C_4$-free
  condition ensures this map is injective on the second component: if both $v_1$
  and $v_2$ are common neighbors of $\{u, w\}$ with $v_1 \neq v_2$, then
  $v_1$-$u$-$v_2$-$w$ forms a $4$-cycle, a contradiction.
\end{proof}

\section{Sperner's Lemma}

In 1912, Luitzen Brouwer published his famous fixed point theorem:
\begin{quote}
\emph{Every continuous function $f : B^n \to B^n$ of an $n$-dimensional ball to
itself has a fixed point (a point $x \in B^n$ with $f(x) = x$).}
\end{quote}
For dimension $1$, that is for an interval, this follows easily from the intermediate
value theorem, but for higher dimensions Brouwer's proof needed some sophisticated
machinery. It was therefore quite a surprise when in 1928 young Emanuel Sperner
produced a simple combinatorial result from which both Brouwer's fixed point theorem
and the invariance of the dimension under continuous bijective maps could be deduced.
And what's more, Sperner's ingenious lemma is matched by an equally beautiful proof
--- it is just double counting.

We discuss Sperner's lemma, and Brouwer's theorem as a consequence, for
the first interesting case, that of dimension $n = 2$.

\begin{lemma}[Sperner's Lemma]
  \label{sperner}
  \lean{chapter28.sperner_lemma}
  Suppose that some ``big'' triangle with vertices $V_1, V_2, V_3$ is triangulated
  (that is, decomposed into a finite number of ``small'' triangles that fit together edge-by-edge).
  Assume that the vertices in the triangulation get ``colors'' from the set $\{1, 2, 3\}$ such that
  $V_i$ receives the color $i$ (for each $i$), and only the colors $i$ and $j$ are used for vertices
  along the edge from $V_i$ to $V_j$ (for $i \neq j$), while the interior vertices are colored
  arbitrarily with $1$, $2$, or $3$.
  Then in the triangulation there must be a small ``tricolored'' triangle, which has all three
  different vertex colors.
\end{lemma}
\begin{proof}
  \leanok
  \uses{handshaking}
  We will prove a stronger statement: The number of tricolored triangles is not only nonzero,
  it is always \emph{odd}.

  Consider the dual graph to the triangulation, but don't take all its edges
  --- only those which cross an edge that has endvertices with the (different)
  colors $1$ and $2$. Thus we get a ``partial dual graph'' which has degree $1$ at all
  vertices that correspond to tricolored triangles, degree $2$ for all triangles in
  which the two colors $1$ and $2$ appear, and degree $0$ for triangles that do not
  have both colors $1$ and $2$. Thus only the tricolored triangles correspond to
  vertices of odd degree (of degree $1$).

  However, the vertex of the dual graph which corresponds to the outside of
  the triangulation has odd degree: in fact, along the big edge from $V_1$ to $V_2$,
  there is an odd number of changes between $1$ and $2$. Thus an odd number
  of edges of the partial dual graph crosses this big edge, while the other big
  edges cannot have both $1$ and $2$ occurring as colors.

  Now since the number of odd-degree vertices in any finite graph is even (by
  equation~(4)), we find that the number of small triangles with three different
  colors (corresponding to odd inside vertices of our dual graph) is odd.
\end{proof}

With this lemma, it is easy to derive Brouwer's theorem.

\begin{theorem}[Brouwer's Fixed Point Theorem (for $n = 2$)]
  \label{brouwer}
  \lean{chapter28.brouwer_fixed_point_2d}
  Every continuous function $f : B^2 \to B^2$ of a $2$-dimensional ball to
  itself has a fixed point (a point $x \in B^2$ with $f(x) = x$).
\end{theorem}
\begin{proof}
  \leanok
  \uses{sperner}
  Let $\Delta$ be the triangle in $\mathbb{R}^3$ with vertices $e_1 = (1,0,0)$,
  $e_2 = (0,1,0)$, and $e_3 = (0,0,1)$. It suffices to prove that every continuous
  map $f : \Delta \to \Delta$ has a fixed point, since $\Delta$ is homeomorphic to $B^2$.

  We use $\delta(T)$ to denote the maximal length of an edge in a triangulation $T$.
  One can easily construct an infinite sequence of triangulations $T_1, T_2, \dots$
  of $\Delta$ such that the sequence of maximal diameters $\delta(T_k)$ converges to $0$
  (for example by iterated barycentric subdivision).

  For each of these triangulations, we define a $3$-coloring of their vertices $v$
  by setting $\lambda(v) := \min\{i : f(v)_i < v_i\}$, that is, $\lambda(v)$ is the smallest
  index $i$ such that the $i$-th coordinate of $f(v) - v$ is negative. If this smallest index $i$
  does not exist, then we have found a fixed point and are done: To see this, note that every
  $v \in \Delta$ lies in the plane $x_1 + x_2 + x_3 = 1$, hence $\sum_i v_i = 1$. So if
  $f(v) \neq v$, then at least one of the coordinates of $f(v) - v$ must be negative (and at
  least one must be positive).

  Let us check that this coloring satisfies the assumptions of Sperner's lemma.
  First, the vertex $e_i$ must receive color $i$, since the only possible negative
  component of $f(e_i) - e_i$ is the $i$-th component. Moreover, if $v$ lies on the
  edge opposite to $e_i$, then $v_i = 0$, so the $i$-th component of $f(v) - v$ cannot
  be negative, and hence $v$ does not get the color $i$.

  Sperner's lemma now tells us that in each triangulation $T_k$ there is a tricolored
  triangle $\{v^{k:1}, v^{k:2}, v^{k:3}\}$ with $\lambda(v^{k:i}) = i$.
  Since the simplex $\Delta$ is compact, some subsequence of $(v^{k:1})_{k \geq 1}$ has a
  limit point~$v$. The distances of $v^{k:2}$ and $v^{k:3}$ from $v^{k:1}$ are at most the
  mesh length $\delta(T_k) \to 0$, so all three sequences converge to the same point~$v$.

  But where is $f(v)$? We know that the first coordinate of $f(v^{k:1})$ is smaller
  than that of $v^{k:1}$ for all $k$. Now since $f$ is continuous, we derive that the
  first coordinate of $f(v)$ is smaller or equal to that of $v$. The same reasoning
  works for the second and third coordinates. Thus none of the coordinates
  of $f(v) - v$ is positive --- and we have already seen that this contradicts
  the assumption $f(v) \neq v$.
\end{proof}
